\section{Quantitative convergence of the arithmetic Gaussian resolvents}
\label{sec:gaussian-resolvent-convergence}

Retain the actual measure and every coordinate of
Section~\ref{sec:kernel-resolvent}:
\[
 s=\tfrac12+it,\qquad
 d\nu_Z(t)=\left|\frac{g(1/2+it)}{h(1/2+it)}\right|^2\frac{dt}{2\pi},
 \qquad E=\mathbb C[s]/(h),\quad A=M_s,\quad c=j_h1.
\]
The polynomial \(h\) contains the complete multiplicities of its
selected actual zeros.  The quotient \(g/h\) is therefore entire,
including at the selected critical zeros.  Use the initial Gaussian
observation of (KR.5) and (KR.8),
\[
 m_n(z)=\sum_{j=1}^n\frac{\omega_{n,j}}{z-\lambda_{n,j}},\qquad
 m(z)=\int\frac{d\nu_Z(t)}{z-s},\qquad
 \mathcal E_n(z)=m(z)-m_n(z),\qquad
 \delta(z)=|\operatorname{Re}z-\tfrac12|.
                                                        \tag{GC.1}
\]
The Gaussian weights are positive and the distinct nodes lie on the
original Mellin line, as proved in Section~\ref{sec:kernel-resolvent}.
Its exact quadrature (KR.8) and exact error (KR.16) are the inputs
below.  No asymptotic for \(q_n\) or \(\kappa_n\) is assumed.

For each fixed \(0<a<\pi/2\), the number
\[
 E_a=\max\left\{\int e^{at}\,d\nu_Z(t),
                       \int e^{-at}\,d\nu_Z(t)\right\}
                                                        \tag{GC.2}
\]
is finite.  Here is the connection to the proved arithmetic decay.
The actual theta calculation in Section~\ref{sec:jet-topology} gives
\(|g(1/2+it)|\leq C_b e^{-b|t|}\) for every \(0<b<\pi/4\).
Choose \(a/2<b<\pi/4\).  Outside a sufficiently large compact interval,
the monic polynomial \(h\) of degree \(d\) satisfies
\(|h(1/2+it)|\geq c_h(1+|t|)^d\).  Thus the density multiplied
by \(e^{a|t|}\) is bounded there by a constant times
\((1+|t|)^{-2d}e^{-(2b-a)|t|}\).  On the compact interval the entire
quotient is bounded.  Both parts are integrable, proving (GC.2) with
the original factor \(dt/(2\pi)\) retained.

\subsection{The Gaussian rule inherits the original exponential moments}

\begin{lemma}[Exponential moment domination]
For every real \(v\) with \(|v|\leq a\), write
\(\lambda_{n,j}=1/2+it_{n,j}\).  Then
\[
 \sum_j\omega_{n,j}e^{vt_{n,j}}
       \leq\int e^{vt}\,d\nu_Z(t)\leq E_a.              \tag{GC.3}
\]
\end{lemma}

\begin{proof}
Let \(H\) be the real Hermite interpolation polynomial of degree at
most \(2n-1\) whose values and first derivatives at the distinct real
nodes \(t_{n,j}\) equal those of \(f(t)=e^{vt}\).  Its existence and
uniqueness follow from the invertibility of the interpolation map:
a polynomial in its kernel has a double zero at each of the \(n\)
nodes, hence is zero in degree below \(2n\); domain and target both
have dimension \(2n\).
For a real \(t\) outside the nodes the interpolation remainder is
\[
 f(t)-H(t)=\frac{f^{(2n)}(\xi)}{(2n)!}
                    \prod_j(t-t_{n,j})^2\geq0           \tag{GC.4}
\]
for some real \(\xi\) between \(t\) and the nodes.  To prove the
remainder, subtract \(C\prod_j(x-t_{n,j})^2\) from \(f(x)-H(x)\),
choosing \(C\) to produce an additional zero at \(x=t\).
The resulting function has the \(n\) double zeros and the additional
zero.  Repeated Rolle's theorem gives a zero of its \(2n\)-th
derivative; solving at that zero gives (GC.4).  The last sign follows
from \(f^{(2n)}(x)=v^{2n}e^{vx}\geq0\).  At a node the remainder is
zero.

Integrate the pointwise inequality.  The exponential and every
polynomial are integrable by (GC.2).  Since \(H\) is a polynomial
in \(t=(s-1/2)/i\) of degree at most \(2n-1\), (KR.8) yields
\[
 \int H(t)\,d\nu_Z(t)
  =\sum_j\omega_{n,j}H(t_{n,j})
  =\sum_j\omega_{n,j}e^{vt_{n,j}}.
\]
This proves the first inequality in (GC.3).  Pointwise convexity of
\(e^{vt}\) in \(v\in[-a,a]\), followed by integration, proves the
second.  No bound on the largest quadrature node has been inserted.
\end{proof}

\subsection{Exact moments give a quantitative analytic-strip estimate}

Define
\[
 \Phi(u)=\int e^{iut}\,d\nu_Z(t),\qquad
 \Phi_n(u)=\sum_j\omega_{n,j}e^{iut_{n,j}},\qquad
 \Delta_n(u)=\Phi(u)-\Phi_n(u).
                                                        \tag{GC.5}
\]
These functions are holomorphic on \(|\operatorname{Im}u|<a\).
Indeed, on each smaller closed strip, every differentiated integrand
is dominated by a constant times \(e^{at}+e^{-at}\), because
\(|t|^k e^{b|t|}\leq C_{k,b,a}(e^{at}+e^{-at})\) when \(b<a\).
Dominated convergence proves all the derivatives of \(\Phi\); the
finite sum is entire.  Equation (GC.3), applied at
\(v=-\operatorname{Im}u\), gives
\[
 |\Delta_n(u)|\leq2E_a\quad(|\operatorname{Im}u|<a),
 \qquad \Delta_n^{(k)}(0)=0\quad(0\leq k<2n).            \tag{GC.6}
\]
The derivative identities are exactly the moments in (KR.8).

\begin{lemma}[The complete strip bound]
For \(|\operatorname{Im}u|<a\),
\[
 |\Delta_n(u)|\leq2E_a
   \left|\tanh\!\left(\frac{\pi u}{4a}\right)\right|^{2n}.
                                                        \tag{GC.7}
\]
\end{lemma}

\begin{proof}
The explicit holomorphic coordinate maps
\[
 w=\tanh\!\left(\frac{\pi u}{4a}\right),\qquad
 u=\frac{2a}{\pi}\operatorname{Log}\frac{1+w}{1-w}
\]
are inverse bijections between the strip and the unit disk.  For
\(|w|<1\), the fraction lies in the right half-plane, so its
principal logarithm has imaginary part in \((-\pi/2,\pi/2)\),
which proves the asserted strip range.  Compose \(\Delta_n/(2E_a)\)
with this inverse.  By (GC.6), the resulting disk function is bounded
by one and has a zero of order at least \(2n\) at zero.  Divide by
\(w^{2n}\), filling in its removable value at zero.  On each circle
\(|w|=r<1\), its modulus is at most \(r^{-2n}\); the maximum
principle gives this bound inside the circle.  Letting \(r\) increase
to one proves the required bound by \(|w|^{2n}\), hence (GC.7).
\end{proof}

\begin{theorem}[Direct off-line convergence with an explicit rate]
\label{thm:gaussian-resolvent-rate}
For every \(0<a<\pi/2\), every \(n\geq1\), and every
\(z\) off the original Mellin line,
\[
 \boxed{\displaystyle
 |m(z)-m_n(z)|\leq\frac{4aE_a}{\pi}
     \Gamma\!\left(\frac{2a\delta(z)}{\pi}\right)
                 (4n)^{-2a\delta(z)/\pi}.}              \tag{GC.8}
\]
On a compact set \(K\) disjoint from that line the same bound holds
uniformly after replacing \(\delta(z)\) everywhere by
\(\delta_K=\min_{z\in K}\delta(z)>0\).
\end{theorem}

\begin{proof}
Put \(w=z-1/2\).  On \(\operatorname{Re}w>0\), the elementary
Laplace identity and Fubini's theorem give
\[
 \mathcal E_n(z)=\int_0^\infty e^{-wu}\Delta_n(u)\,du.
                                                        \tag{GC.9}
\]
For example the absolute integral for the original measure before
subtraction is at most \(\nu_Z(\mathbb R)/\operatorname{Re}w\),
so the interchange is justified.  On \(\operatorname{Re}w<0\),
the exact orientation is
\[
 \mathcal E_n(z)=-\int_0^\infty e^{wu}\Delta_n(-u)\,du.
                                                        \tag{GC.10}
\]
These follow respectively from integrating \(e^{-(w-it)u}\) and
\(-e^{(w-it)u}\).  Set \(b=\pi/(4a)\).  Equation (GC.7) gives
\[
 |\mathcal E_n(z)|\leq2E_a I_n(\delta(z)),\qquad
 I_n(\delta)=\int_0^\infty e^{-\delta u}\tanh(bu)^{2n}\,du.
                                                        \tag{GC.11}
\]
For \(0<x<1\), direct integration gives
\[
 \log\frac{1-x}{1+x}=-2\int_0^x\frac{dv}{1-v^2}\leq-2x.
\]
Taking \(x=e^{-2bu}\) yields
\(\tanh(bu)^{2n}\leq e^{-4ne^{-2bu}}\).  With
\(\beta=\delta/(2b)=2a\delta/\pi\), the same variable change gives
\[
 \begin{split}
 I_n(\delta)
 &\leq\frac1{2b}\int_0^1x^{\beta-1}e^{-4nx}\,dx\\
 &\leq\frac1{2b}(4n)^{-\beta}\Gamma(\beta),\qquad
 \Gamma(\beta)=\int_0^\infty y^{\beta-1}e^{-y}\,dy.
 \end{split}                                           \tag{GC.12}
\]
Substitution in (GC.11) proves (GC.8).  On the compact set one
uses \(e^{-\delta(z)u}\leq e^{-\delta_Ku}\) in (GC.11) and
then (GC.12) at \(\delta_K\); no monotonicity of the gamma factor
is assumed.  This proves locally uniform convergence on both
components of the off-line domain.
\end{proof}

\subsection{Every confluent derivative and the original finite algebra}

Write \(F^{[k]}=F^{(k)}/k!\).  The exact error (KR.16), with its
vertical-line sign, gives for all \(k\geq0\)
\[
 \boxed{\displaystyle
 \mathcal E_n^{[k]}(z)=(-1)^n\sum_{j=0}^k
   (q_n^{-2})^{[k-j]}(z)(-1)^j
          \int\frac{|q_n(s)|^2}{(z-s)^{j+1}}\,d\nu_Z(t).}
                                                        \tag{GC.13}
\]
Every derivative under this integral is justified on a compact
off-line neighborhood by its positive distance from the line and
the finite measure \(|q_n|^2\nu_Z\).  The product rule for Taylor
coefficients has no binomial factors; those are precisely cancelled
by the convention \(F^{[k]}=F^{(k)}/k!\).  For completeness its
polynomial factor is explicitly
\[
 (q_n^{-2})^{[r]}(z)=\frac{(-1)^r}{q_n(z)^2}
   \sum_{r_1+\cdots+r_n=r}
        \prod_{j=1}^n\frac{r_j+1}{(z-\lambda_{n,j})^{r_j}}.
                                                        \tag{GC.14}
\]
This follows by multiplying the Taylor series of the \(n\) original
factors \((z-\lambda_{n,j})^{-2}\), retaining every cross term.

Differentiating (GC.9)--(GC.10) gives the independent estimate
\[
 |\mathcal E_n^{[k]}(z)|\leq\frac{2E_a}{k!}
   \int_0^\infty u^k e^{-\delta(z)u}\tanh(bu)^{2n}\,du.  \tag{GC.15}
\]
The integrable majorant \(u^k e^{-\delta_Ku}\) on a smaller
off-line compact neighborhood justifies this differentiation.
If \(0<\sigma<\delta\), the exponential series gives
\(u^k e^{-(\delta-\sigma)u}\leq k!/(\delta-\sigma)^k\).
Equations (GC.11)--(GC.12) prove
\[
 \boxed{\displaystyle
 |\mathcal E_n^{[k]}(z)|\leq
 \frac{4aE_a}{\pi(\delta-\sigma)^k}
       \Gamma\!\left(\frac{2a\sigma}{\pi}\right)
                  (4n)^{-2a\sigma/\pi},
 \quad 0<\sigma<\delta\leq\delta(z).}                   \tag{GC.16}
\]
Taking \(\delta=\delta_K\) makes this a simultaneous compact-uniform
bound on every fixed finite list of derivatives.

Here is the complete algebra transport, including the part of the
original packet outside this resolvent domain.  Write
\[
 h(s)=\prod_{\rho\in Z}(s-\rho)^{d_\rho},\qquad
 H_\rho(s)=\frac{h(s)}{(s-\rho)^{d_\rho}},\qquad
 U_\rho(s)=\sum_{k=0}^{d_\rho-1}
          (H_\rho^{-1})^{[k]}(\rho)(s-\rho)^k,
\]
and define the actual polynomial remainder
\[
 e_\rho=j_h(H_\rho U_\rho)\in E.                        \tag{GC.17}
\]
Its remainder is one modulo \((s-\rho)^{d_\rho}\) and zero
modulo every other factor of \(h\).  The distinct factors are
coprime, and a polynomial divisible by each is divisible by their
product: successive use of the Bezout identity proves this last
statement.  It follows directly that
\[
 e_\rho e_\eta=0\ (\rho\ne\eta),\qquad
 e_\rho^2=e_\rho,\qquad \sum_{\rho\in Z}e_\rho=1.        \tag{GC.18}
\]
Restriction of a polynomial to its truncated Taylor coefficients
therefore gives the exact algebra isomorphism
\[
 E\longrightarrow\bigoplus_{\rho\in Z}
        \mathbb C[\epsilon_\rho]/(\epsilon_\rho^{d_\rho}),
 \qquad
 j_hP\longmapsto
    \left(\sum_{k<d_\rho}P^{[k]}(\rho)\epsilon_\rho^k\right)_\rho.
                                                        \tag{GC.19}
\]
Its inverse sends the coordinate \(\epsilon_\rho^k\) to
\(e_\rho(s-\rho)^k\); composition is the identity by the just
proved congruences.  Multiplication is preserved because the Taylor
coefficient of a product is the exact convolution of the two
coefficient lists, truncated only by the displayed ideal.

Let \(Z_{\rm off}=\{\rho\in Z:\operatorname{Re}\rho\ne1/2\}\)
and \(Z_{\rm line}=Z\setminus Z_{\rm off}\).  Define
\[
 e_{\rm off}=\sum_{\rho\in Z_{\rm off}}e_\rho,\qquad
 e_{\rm line}=1-e_{\rm off},\qquad
 E_{\rm off}=e_{\rm off}E,\quad E_{\rm line}=e_{\rm line}E.
                                                        \tag{GC.20}
\]
Then \(E=E_{\rm off}\oplus E_{\rm line}\) as algebras with
their respective idempotent units.  The projection
\(u\mapsto e_{\rm off}u\) has exactly kernel \(E_{\rm line}\)
and commutes with the original multiplication operator \(A\).
Under (GC.19), its restriction at every retained centre is
\(A_\rho=\rho I+N_\rho\), where multiplication by
\(\epsilon_\rho\) gives \(N_\rho^{d_\rho}=0\) and retains all
lower powers.  This constructs the precise map from the full
packet to the off-line resolvent domain, together with its full
kernel and the operator on that kernel.  An empty off-line set
gives the zero summand, with the full original packet in the kernel.

For a function \(f\) holomorphic near \(Z_{\rm off}\), define
\[
 f(A_{\rm off})=
  \sum_{\rho\in Z_{\rm off}}\sum_{k=0}^{d_\rho-1}
     f^{[k]}(\rho)(A-\rho I)^k e_\rho(A)\big|_{E_{\rm off}}.
                                                        \tag{GC.21}
\]
Here \(e_\rho(A)\) means multiplication by the explicitly
constructed remainder (GC.17).  Equation (GC.19) proves that
(GC.21) is a unital algebra homomorphism on the off-line summand.
No diagonalizability of \(A\) has entered this construction.

In particular, finite geometric expansion in each nilpotent block
proves
\[
 \begin{split}
 m(A_{\rm off})
   &=\int(A_{\rm off}-sI)^{-1}\,d\nu_Z(t),\\
 m_n(A_{\rm off})
   &=\sum_j\omega_{n,j}(A_{\rm off}-\lambda_{n,j}I)^{-1},\\
 \mathcal E_n(A_{\rm off})
   &=\sum_{\rho\in Z_{\rm off}}\sum_{k<d_\rho}
     \mathcal E_n^{[k]}(\rho)(A-\rho I)^ke_\rho(A)
                              \big|_{E_{\rm off}}.
 \end{split}                                           \tag{GC.22}
\]
For the first equality the expansion is
\[
 (A-sI)^{-1}e_\rho(A)=
  \sum_{k<d_\rho}\frac{(-1)^k(A-\rho I)^ke_\rho(A)}
                             {(\rho-s)^{k+1}}.
\]
Each coefficient is integrable since \(\delta(\rho)>0\), and
its integral is exactly \(m^{[k]}(\rho)\).  This proves the
integral identity, rather than assigning a formal meaning to it.

For a concrete bound keep any fixed positive Hermitian form
\(G_*\) on the original coefficient space and restrict it to
\(E_{\rm off}\).  When that summand is nonzero, put
\[
 \delta_* =\min_{\rho\in Z_{\rm off}}\delta(\rho)>0,
 \qquad
 C_{\rho,k}=\left\|(A-\rho I)^ke_\rho(A)
                      \big|_{E_{\rm off}}\right\|_{G_*}.
\]
The triangle inequality and (GC.16) give, for every
\(0<\sigma<\delta_*\),
\[
 \boxed{\displaystyle
 \|m(A_{\rm off})-m_n(A_{\rm off})\|_{G_*}
 \leq\frac{4aE_a}{\pi}
       \Gamma\!\left(\frac{2a\sigma}{\pi}\right)
       (4n)^{-2a\sigma/\pi}
       \sum_{\rho\in Z_{\rm off}}\sum_{k<d_\rho}
             \frac{C_{\rho,k}}{(\delta_*-\sigma)^k}.}     \tag{GC.23}
\]
Thus the full initial resolvent converges in the original finite
coefficient space, with every nilpotent contribution controlled.
Applying the displayed operators to the retained vector
\(e_{\rm off}c\) gives the corresponding vector statement with
the additional factor \(\|e_{\rm off}c\|_{G_*}\).

The polynomial kernel (KR.9) continues to act on all of \(E\),
including \(E_{\rm line}\) and collisions with quadrature nodes.
The off-line estimates above concern the initial observation
\(m_n\).  Its exact map to the terminal observation \(F_n\) is
the block and tail comparison (KR.14)--(KR.20), with the factors
\(q_n,\kappa_n,a_n\) and the tail denominator still present.
The estimate (GC.23) may be inserted into those identities; it
does not remove their arithmetic factors or the critical-line
summand of (GC.20).
